\documentclass[11pt,a4paper]{article}
\usepackage[T1]{fontenc}
\usepackage{lmodern,amsmath,amssymb,amsthm,microtype}
\usepackage[margin=26mm]{geometry}
\usepackage[colorlinks=true,linkcolor=blue,urlcolor=blue]{hyperref}
\newtheorem{theorem}{Theorem}[section]
\newtheorem{corollary}[theorem]{Corollary}
\newtheorem{lemma}[theorem]{Lemma}
\newcommand{\Q}{\mathbb Q}\newcommand{\Z}{\mathbb Z}
\title{A denominator-support obstruction for quotient-height detectors}
\author{}\date{17 November 2026}
\begin{document}\maketitle

\begin{abstract}
If every positive odd prime of a detector weight must occur in the
denominator of one elliptic point, the logarithmic radical of that
weight is bounded by the height of the point. This is incompatible
with a positive radical lower bound controlled by a height coordinate
growing quadratically in an independent index ratio. Applied to the
pole-cancelling coordinate $X^2x([n]H)-x(H)$ and the sixteen-shift
derivative weight, the obstruction bounds $n$ for every fixed source.
Thus that combined source cannot be complete for all integer inputs.
The argument extends to every fixed rational height coordinate that
is generically nonconstant in the second point and satisfies the
stated uniform height estimates. Siegel's theorem controls the
constant specializations. No assertion is made about other
two-point sources.
\end{abstract}

\section{A general denominator bound}

Fix an elliptic curve $E/\Q$ in an integral Weierstrass model and
a nontorsion point $P\in E(\Q)$. For $H=[h]P$, $h\ne0$, write
\[
                    x(H)=A_H/B_H^2,
             \qquad A_H\in\Z,\quad B_H\in\Z_{>0},
             \quad\gcd(A_H,B_H)=1.
\]
The canonical-height comparison \cite[Chapter VIII, Section 9]{AEC}
gives a fixed constant $C_E$ such that
\begin{equation}\label{eq:denominator-upper}
                 \log B_H\le\tfrac12h(x(H))
                                  \le\widehat h(H)+C_E.
\end{equation}
Let $S$ contain the bad reduction primes and put
\[
            r_S^+(u)=\prod_{\substack{p\notin S\\v_p(u)>0\text{ odd}}}p
                       \qquad(u\in\Q^\times).
\]

\begin{theorem}\label{thm:general}
Consider a family of guarded rational source solutions carrying
nonzero integer indices $h,n$, a weight $u$, and a rational
height coordinate $M$. Suppose there are fixed constants
$a,\delta>0$ and $C_0,C_1$ such that every source solution satisfies
\begin{align}
 v_p(u)>0\text{ odd},\ p\notin S
       &\quad\Longrightarrow\quad H\text{ reduces to }O\pmod p,
                                                        \label{eq:formal}\\
 \log r_S^+(u)&\ge\delta h(M)-C_0,                       \label{eq:radical}\\
 h(M)&\ge a(n^2-1)\widehat h(H)-C_1.                    \label{eq:carrier}
\end{align}
Then $|n|$ is bounded by a constant depending only on the fixed
source. If $3a\delta>1$, the indices $h$ are also bounded on
the subfamily $|n|\ge2$.
\end{theorem}
\begin{proof}
At a good prime, reduction of $H$ to the identity is equivalent
to $p\mid B_H$. Thus \eqref{eq:formal} gives
\[
          r_S^+(u)\mid\operatorname{rad}(B_H),\qquad
          \log r_S^+(u)\le\log B_H\le\widehat h(H)+C_E.
\]
Combining this with \eqref{eq:radical} and \eqref{eq:carrier}
yields
\begin{equation}\label{eq:obstruction}
       \bigl(a\delta(n^2-1)-1\bigr)\widehat h(H)
                                  \le C_E+C_0+\delta C_1.
\end{equation}
Since $\widehat h(H)=h^2\widehat h(P)\ge\widehat h(P)>0$,
this bounds $n$. If $|n|\ge2$ and $3a\delta>1$, the coefficient
on the left is at least $3a\delta-1>0$, so it bounds $h$ too.
\end{proof}

No lower estimate for the size of an elliptic denominator is
used here. The upper bound \eqref{eq:denominator-upper} is
enough. The support condition \eqref{eq:formal} can be replaced
by the weaker assertion that every such prime divides $B_H$.

\section{The pole-cancelling coordinate and derivative weight}

Take $E_d:y^2=x^3-d^2x$, $U=[n]H$, and
\[
 \beta=X^2+B,\qquad Y=X^2x(U)-x(H),\qquad
 E_h=h(1:\beta:Y),\qquad M=\beta^4+Y^4,
\]
where $X\ne0$, $B\in\Z_{>0}$ is fixed, and all indices are
nonzero. Set $a_0=M^3$, $t_0=M^6-z^2\ne0$, and
\[
 F(a_0,z)=\prod_{c\in\{-8,\ldots,-1,1,\ldots,8\}}
                    ((a_0+c)^2-z^2),\qquad
 u=-\frac{32a_0F(a_0,z)}{F_{a_0}(a_0,z)}.
\]
Use the guarded one-slot balanced detector containing
$C_1=1/x(H)$, together with the intrinsic pole condition for
$7b$, where
\[
 D_i=\frac{C_i^2}{1+uC_i^2},\qquad
 W=\prod_i(1+D_i^2/u),\qquad
 b=\frac{4uW^2}{(1+u)^2}.
\]
Any additional tests may be retained.

\begin{corollary}\label{cor:application}
Under ordinary rational $abc$, every fixed source of this form
has bounded $|n|$. Hence the source is not complete for the
intended inputs $X=n\in\Z\setminus\{0\}$.
For the retained constants, both $n$ and $h$ are bounded on
$|n|\ge2$. If all other fixed data are held independent of
$B$, then sufficiently large $B$ eliminates that entire
subfamily.
\end{corollary}
\begin{proof}
Include two, seven and the fixed exceptional primes in $S$.
At a positive odd valuation $\nu=v_p(u)$, every factor of
$W$ has valuation at most zero. If $H$ is not formal, then
$v_p(C_1)\le0$, so the corresponding factor has valuation
at most $-\nu$. Consequently
\[
                       v_p(b)=\nu+2v_p(W)<0
\]
is odd, contradicting the intrinsic pole condition at $p\ne7$.
This proves \eqref{eq:formal}.

The derivative-height and signed mixed-power estimates
\cite{Derivative,Height} give
\[
 \log r_S^+(u)\ge\gamma h(t_0)-C h(M)-O(1),\qquad
 h(M)\le Qh(t_0)+O(1),
\]
with
\[
 \gamma=\frac{997}{200},\qquad C=\frac{401}{100},\qquad
 Q=\frac{1002}{995}.
\]
Thus \eqref{eq:radical} holds with
\[
               \delta=\frac\gamma Q-C=\frac{188411}{200400}>0.
\]
The positive-power and pole-cancelling height bounds
\cite{Height,Cancellation} give
\[
 h(M)\ge\alpha E_h-O(1),\qquad
 E_h\ge(n^2-1)\widehat h(H)-O(1),\qquad
 \alpha=\frac{997}{1001}.
\]
Apply Theorem~\ref{thm:general} with $a=\alpha$. The exact
coefficient in \eqref{eq:obstruction} is
\[
             a\delta=\frac{187845767}{200600400}>0,
 \qquad 3a\delta-1=\frac{120978967}{66866800}>0.
\]
This bounds both indices on $|n|\ge2$. Equations
\eqref{eq:denominator-upper} and \eqref{eq:radical} then
bound $h(M)$. The lower estimate $h(M)\ge\alpha E_h-O(1)$
and $E_h\ge h(\beta)\ge\log B$ contradict arbitrarily large
$B$ when the constants are uniform in that shift.
\end{proof}

The obstruction concerns a detector whose entire positive odd
support must be carried by the denominator of $H$, together
with the stated lower bound for its height coordinate. It does
not exclude two-point constructions using different support
conditions or different ways to identify the input. The local
cancellation of $Y$ at $H=O$ remains valid; it does not enlarge
the available denominator support.

\section{Fixed rational height coordinates}

The same obstruction applies to a fixed rational replacement for the
pole-cancelling coordinate. Fix a rational function
\[
 R\in\Q(\mathbb A^1_X\times E_H\times E_U)
\]
that is generically nonconstant in $U$. Work on a fixed affine chart
where the retained parameter denominators are nonzero. Reducing
powers of $y(U)$ modulo its curve equation and clearing a common
parameter denominator gives
\begin{equation}\label{eq:carrier-form}
 R=\frac{A(Z)+B_1(Z)y(U)}{C(Z)+D(Z)y(U)},\qquad Z=x(U),
\end{equation}
where the coefficients are fixed polynomials in $X,x(H),y(H)$.
The displayed denominator is required to be nonzero at a source
solution. Write $p,q$ for the coefficient vectors of $A+B_1y$
and $C+Dy$, in the same ordered basis of monomials $Z^i,Z^iy$.
Put
\[
                   \Gamma=\sum_{i<j}(p_iq_j-p_jq_i)^2.
\]
At rational parameters, $\Gamma\ne0$ is equivalent to the
specialization of $R$ being nonconstant in $U$.

\begin{lemma}\label{lem:carrier-height}
There are fixed constants $K_X,K_H\ge0$ such that, at every
guarded rational specialization with $\Gamma\ne0$,
\[
 h\bigl(R(X,H,U)\bigr)
       \ge\widehat h(U)-K_H\widehat h(H)-K_Xh(X)-O(1).
\]
If the coefficient polynomials in \eqref{eq:carrier-form} have
degrees at most $d_X,d_x,d_y$ separately in $X,x(H),y(H)$,
one may take $K_X=d_X$ and $K_H=2d_x+3d_y$.
\end{lemma}
\begin{proof}
Let $r=R(X,H,U)$. The polynomial
\[
 L_r(Z)=(A(Z)-rC(Z))^2
              -(Z^3-d^2Z)(B_1(Z)-rD(Z))^2
\]
vanishes at $x(U)$. It is not the zero polynomial: the cubic
$Z^3-d^2Z$ is not a square in $\Q(Z)$, so the contrary would
force both $A-rC$ and $B_1-rD$ to vanish, making $R$ constant.
The rational-root theorem, after removing any zero-power factor,
and the fixed polynomial coefficient bounds give
\[
 h(x(U))\le h(L_r)
              \le2h(r)+2h(p:q)+O(1).
\]
Here $h(L_r)$ denotes projective coefficient height. Polynomial
evaluation and the elliptic height comparisons give
\[
 h(p:q)\le d_Xh(X)+d_xh(x(H))+d_yh(y(H))+O(1)
 \le d_Xh(X)+(2d_x+3d_y)\widehat h(H)+O(1).
\]
Finally $h(x(U))\ge2\widehat h(U)-O(1)$, which proves the claim.
\end{proof}

\begin{theorem}\label{thm:fixed-carrier}
Fix $B\ge1$ and a rational function $R$ as above. Consider guarded
rational source solutions with $H=[h]P$, $U=[n]H$,
$h,n\in\Z\setminus\{0\}$, and put
\[
             E_R=h(1:X^2+B:R(X,H,U)).
\]
Suppose fixed positive constants $\alpha,\delta$ give
\[
 r_S^+(u)\mid\operatorname{rad}(B_H),\qquad
 \log r_S^+(u)\ge\delta h(M)-O(1),\qquad
 h(M)\ge\alpha E_R-O(1).
\]
On the subfamily $\Gamma\ne0$, the index $n$ is bounded.
Without imposing $\Gamma\ne0$, only finitely many integers $n$
can have intended multiplication witnesses with $X=n$.
\end{theorem}
\begin{proof}
The hypotheses give
\begin{equation}\label{eq:general-carrier-budget}
                       \alpha\delta E_R\le\widehat h(H)+O(1).
\end{equation}
Since $B\ge1$, one has $E_R\ge2h(X)$ and $E_R\ge h(R)$.
At a nonconstant specialization Lemma~\ref{lem:carrier-height}
therefore gives
\[
 E_R\ge\frac{2}{K_X+2}(n^2-K_H)\widehat h(H)-O(1).
\]
Combining with \eqref{eq:general-carrier-budget} yields
\[
 \left(\frac{2\alpha\delta}{K_X+2}(n^2-K_H)-1\right)
                         \widehat h(H)\le O(1),
\]
so $n$ is bounded, because $\widehat h(H)\ge\widehat h(P)>0$.

It remains to consider $\Gamma=0$. At rational parameters this
is the common zero locus of the coefficient minors. That locus
is a proper algebraic subset of $\mathbb A^1_X\times E_H$,
since $R$ is generically nonconstant in $U$. Take the closure of each of its curve
components in $\mathbb P^1\times E$. Zero-dimensional and
singular exceptions contribute finitely many pairs $(X,H)$.
A geometric component not defined over $\Q$ has only finitely
many rational points, since these lie in its intersection with
a distinct Galois conjugate. It remains to consider a
geometrically integral component defined over $\Q$ and its
smooth projective normalization $C_0$.

If the $X$-projection is constant, it contributes at most one
rational input. If the $H$-projection is constant, then
\eqref{eq:general-carrier-budget} and $E_R\ge2h(X)$ bound
$h(X)$, hence give finitely many rational inputs. In the remaining
case both projections are nonconstant. The map $C_0\to E$
forces $g(C_0)\ge1$ by Riemann--Hurwitz. Siegel's theorem in
rational-function form \cite[Theorem 1.1]{Levin} then gives
finitely many rational points of $C_0$ at which the nonconstant
function $X$ takes an integer value. This proves finiteness of
the intended integer inputs on the degeneration locus as well.
\end{proof}

The second conclusion concerns witnesses with $X=n$ and
$U=[n]H$. It is not a statement about the full rational-input
projection of a source that has not independently identified
$X$ with $n$. The theorem also requires one fixed rational
function $R$, generically nonconstant in $U$, and the stated
uniform radical and height estimates. It does not rule out
different support conditions, functions independent of $U$,
or additional free height coordinates.

\begin{thebibliography}{9}
\bibitem{AEC}J.~H.~Silverman, \emph{The Arithmetic of Elliptic Curves},
second edition, Graduate Texts in Mathematics 106, Springer, 2009,
Chapter VIII, Section 9.
\bibitem{Derivative}\emph{An eleven-variable polynomial under ordinary
$abc$}, companion paper, \nolinkurl{derivative_detector11_20261117_v59.tex}.
\bibitem{Height}\emph{Signed mixed-power heights and single-cover
criteria}, companion paper, \nolinkurl{signed_mixed_power_height_v58.tex}.
\bibitem{Cancellation}\emph{A pole-cancelling height coordinate for an
elliptic quotient}, companion paper,
\nolinkurl{pole_cancelling_quotient_height_v59.tex}.
\bibitem{Levin}A.~Levin, \emph{Siegel's theorem and the Shafarevich
conjecture}, Journal de Th\'eorie des Nombres de Bordeaux 24 (2012),
705--727, Theorem 1.1. \href{https://doi.org/10.5802/jtnb.818}{doi:10.5802/jtnb.818}.
\end{thebibliography}
\end{document}
